El Capítulo~\ref{chp:conocimiento_cero} presentó los sistemas de prueba y las dos herramientas con
que el presente trabajo aborda el problema de los puentes descrito en la
Sección~\ref{sec:puentes}. Este capítulo traslada esos fundamentos a un puente concreto, en el que
la cadena de destino acuña un activo solo cuando verifica pruebas de que una transacción que invoca
la función de bloqueo fue incluida en un bloque de la cadena de origen firmado por su comité de
sincronización. Como se precisa en la Sección~\ref{sec:puente_aceptacion}, estas pruebas no
acreditan que dicha llamada se haya ejecutado con éxito ni que el bloque haya alcanzado la
finalidad.

La Sección~\ref{sec:lock_and_mint} describe la arquitectura común del puente, el flujo de una
transferencia y las condiciones bajo las que se autoriza la acuñación. Las secciones siguientes
reseñan zkBridge, de Berkeley RDI, como referencia de diseño, presentan el diseño y la
implementación de las pruebas en Noir y en SP1 y exponen la metodología con que ambas se comparan
y los resultados obtenidos.

\section{Implementación de un puente lock-and-mint}
\label{sec:lock_and_mint}

El puente transfiere el token \texttt{GuaraniToken} desde una cadena de origen hacia una cadena de
destino mediante el esquema \textit{lock-and-mint}. El presente trabajo implementa únicamente el
flujo de ida, de modo que la operación inversa, de quema y liberación, queda fuera de su alcance.
La descripción que sigue corresponde a la documentación y al código de la revisión del 12 de
septiembre de 2026 del repositorio del puente~\cite{zkguaranibridgesp1}.

\subsection{Arquitectura y componentes}
\label{sec:puente_arquitectura}

Conviene distinguir dos variantes del esquema. En el puente de referencia descrito en un trabajo previo de la
misma institución, un \textit{relayer} autorizado observa el bloqueo en la cadena de origen e invoca la
acuñación en la de destino, cuyo contrato receptor comprueba únicamente que la llamada proviene de
ese \textit{relayer} y que el identificador del bloqueo no se procesó antes, por lo que la
corrección de cada acuñación depende por completo de su honestidad~\cite{figueredo2026zkbridge,
guaranibridgevanilla}. En la variante que desarrolla el presente trabajo, el contrato de destino
acuña solo tras verificar dos pruebas, y el \textit{relayer} deja de contar con un permiso
exclusivo, ya que cualquier parte puede presentar esas pruebas y solicitar la acuñación. Las garantías
de una variante no se trasladan a la otra, y lo que sigue se refiere exclusivamente a la segunda.

La arquitectura se compone de los siguientes elementos.
\begin{itemize}
	\item \texttt{GuaraniToken}, un token ERC-20 con símbolo GUA desplegado en ambas cadenas. En la
	de origen se emite con un suministro inicial, y en la de destino se despliega sin suministro y
	solo el contrato receptor recibe el rol que autoriza la acuñación.
	\item \texttt{Sender}, el contrato de bloqueo en la cadena de origen. Su función
	\lstinline[breaklines=true]|lock(recipientL2, amount)| exige un destinatario distinto de la dirección cero, un
	monto positivo y una autorización de gasto suficiente, transfiere los tokens a su custodia,
	registra el bloqueo y emite el evento \texttt{Locked}. Este contrato y el token se conservan
	sin cambios respecto del prototipo de ese trabajo previo, del que deriva el repositorio
	según su historial de versiones~\cite{figueredo2026zkbridge, zkguaranibridgesp1}.
	\item \texttt{BeaconLightClient}, en la cadena de destino, que verifica la prueba de consenso y
	registra como anclado el hash del bloque de ejecución que esta certifica. Su propietario fija
	los parámetros de política que la prueba no puede validar por sí sola, que son el dominio de
	firma, el compromiso del comité de sincronización vigente y la posición del hash del bloque en
	el encabezado de la cadena de baliza.
	\item \texttt{ZKReceiver}, en la cadena de destino, que verifica la prueba de inclusión y acuña
	los tokens. Sus funciones de acuñación son públicas, de modo que cualquier parte que disponga de
	pruebas válidas puede invocarlas.
	\item Un contrato verificador, en el que ambos contratos de destino delegan la comprobación
	criptográfica de las pruebas.
	\item El \textit{relayer}, un proceso en Node.js que se ejecuta fuera de las cadenas, detecta los
	bloqueos, obtiene los datos de la cadena de origen, invoca los programas que generan las pruebas
	y envía la operación a la cadena de destino. Los contratos no le reconocen ningún permiso
	exclusivo ni aceptan como parámetro de la acuñación ningún dato suyo que no esté comprometido
	por una prueba, aunque su disponibilidad condiciona que las transferencias se completen.
	\item Una interfaz web que permite al usuario conectar su billetera, autorizar el gasto, invocar
	el bloqueo y seguir los eventos de bloqueo y de acuñación.
\end{itemize}

Las pruebas establecen dos enunciados enlazados por el hash del bloque de ejecución. La prueba de
consenso establece que el comité de sincronización firmó un encabezado de la cadena de baliza del
que se deriva ese hash, y la prueba de inclusión, que la transacción que invoca \texttt{lock}
pertenece al bloque cuyo encabezado produce el mismo hash. Ninguno de los dos establece que esa
llamada se haya ejecutado con éxito. Los contratos, el token, el \textit{relayer}, la interfaz y
estos dos enunciados son comunes a las implementaciones en Noir y en SP1. La forma en que cada
herramienta expresa y demuestra los enunciados, la codificación de sus valores públicos y el
contrato verificador que corresponde a cada sistema de prueba se describen en las secciones
dedicadas a cada una de ellas.

\begin{figure}[htbp]
	\centering
	\begin{tikzpicture}[
		comp/.style = {draw, rounded corners, align=center, minimum height=0.95cm,
			text width=2.45cm, font=\footnotesize},
		dato/.style = {font=\scriptsize, align=center},
		flecha/.style = {-{Stealth[length=2.2mm]}, thick},
		doble/.style = {{Stealth[length=2.2mm]}-{Stealth[length=2.2mm]}, thick},
		evento/.style = {-{Stealth[length=2.2mm]}, thick, dashed}
	]
		\node[comp] (usuario) at (0, 3.5) {Usuario e\\interfaz web};
		\node[comp] (sender) at (0, 1.1) {\texttt{Sender}};
		\node[comp] (token1) at (0, -1.1) {\texttt{GuaraniToken}};
		\node[comp] (relayer) at (4.6, 1.1) {\textit{Relayer}};
		\node[comp] (programas) at (4.6, 3.5) {Programas de\\consenso e inclusión};
		\node[comp] (datos) at (4.6, -1.1) {Nodos de la\\cadena de origen};
		\node[comp] (receiver) at (9.2, 1.1) {\texttt{ZKReceiver}};
		\node[comp, text width=2.75cm] (light) at (9.2, -1.1) {{\ttfamily\scriptsize BeaconLightClient}};
		\node[comp] (token2) at (13.5, 1.1) {\texttt{GuaraniToken}};
		\node[comp] (verif) at (13.5, -1.1) {Contrato\\verificador};

		\node[draw, dashed, rounded corners, inner sep=6pt, fit=(sender) (token1),
			label={[font=\footnotesize]below:Cadena de origen}] {};
		\node[draw, dashed, rounded corners, inner sep=6pt, fit=(receiver) (light) (token2) (verif),
			label={[font=\footnotesize]below:Cadena de destino}] {};

		\draw[flecha] (usuario) -- node[right, dato] {1. \texttt{approve}, \texttt{lock}} (sender);
		\draw[flecha] (sender) -- node[right, dato] {custodia} (token1);
		\draw[evento] (sender) -- node[above, dato] {\texttt{Locked}} (relayer);
		\draw[doble] (relayer) -- node[right, dato] {2. recibo, bloque,\\comité} (datos);
		\draw[doble] (relayer) -- node[right, dato] {3. entradas\\y pruebas} (programas);
		\draw[flecha] (relayer) -- node[above, dato, pos=0.45] {4. pruebas} (receiver);
		\draw[flecha] (receiver) -- node[above, dato] {5. \texttt{mint}} (token2);
		\draw[flecha] (receiver) -- node[left, dato] {anclaje} (light);
		\draw[flecha] (receiver) -- (verif);
		\draw[flecha] (light) -- (verif);
	\end{tikzpicture}
	\caption{Arquitectura del puente y recorrido de una transferencia de ida. Las líneas
	discontinuas indican eventos y los números, las etapas descritas en el texto.}
	\label{fig:arquitectura_puente}
\end{figure}

\subsection{Flujo operativo de la transferencia}
\label{sec:puente_flujo}

La Figura~\ref{fig:arquitectura_puente} resume la arquitectura y el recorrido de una transferencia
de ida, que se desarrolla en cinco etapas.

\begin{enumerate}
	\item \textbf{Solicitud y bloqueo.} El usuario autoriza al contrato \texttt{Sender} a transferir
	sus tokens e invoca \lstinline|lock(recipientL2, amount)|. Si las comprobaciones del contrato se
	cumplen, los tokens quedan en custodia y se emite el evento \texttt{Locked}. Los datos de la
	llamada, es decir, el selector de la función, el destinatario y el monto, son los que el destino
	utilizará más adelante.
	\item \textbf{Detección y obtención de datos.} El \textit{relayer} detecta el evento, ya sea por
	suscripción o por consulta periódica, y obtiene el recibo de la transacción para identificar el
	bloque que la contiene. Con un nodo de la cadena de baliza obtiene los datos del comité de
	sincronización que firmó ese bloque, y con un nodo de ejecución reconstruye el encabezado del
	bloque y obtiene todas sus transacciones. El \textit{relayer} comprueba que el hash del
	encabezado reconstruido y el de cada transacción coincidan con los que informa el nodo, y
	detiene el proceso si no coinciden. Los datos del evento solo sirven para iniciar el proceso y
	no se utilizan para acuñar.
	\item \textbf{Generación de las pruebas.} El \textit{relayer} invoca el programa de consenso con
	el encabezado de la cadena de baliza, las claves públicas del comité, el vector de
	participación, la firma agregada y la rama de Merkle que conduce al hash del bloque de
	ejecución, y el programa de inclusión con el encabezado de
	ejecución, las transacciones del bloque y la posición de la transacción de bloqueo. Cada
	programa produce una prueba y sus valores públicos.
	\item \textbf{Envío a destino.} Si el bloque todavía no está anclado, el \textit{relayer} envía
	ambas pruebas en una única transacción a \lstinline|mintRemoteWithHeader|, junto con los datos
	de la llamada de bloqueo, que no forman parte de los valores públicos. Si el bloque ya fue
	anclado por otra operación, envía solo la prueba de inclusión a \lstinline|mintRemote|. En
	ambos casos espera después el evento \texttt{Minted} para confirmar la operación.
	\item \textbf{Comprobaciones y acuñación.} Las comprobaciones dependen de la ruta utilizada en
	la etapa anterior. Con \lstinline|mintRemoteWithHeader|, el contrato \texttt{ZKReceiver}
	solicita primero a \texttt{BeaconLightClient} que verifique la prueba de consenso y ancle el bloque que esta
	certifica. Con \lstinline|mintRemote| omite ese paso y se apoya en el anclaje registrado por una
	operación anterior, sin volver a verificar la prueba que lo produjo. En ambos casos verifica
	después la prueba de inclusión, exige que el bloque de la transacción probada figure como anclado
	y, si se cumplen las demás condiciones de la Sección~\ref{sec:puente_aceptacion}, acuña el monto
	a favor del destinatario y emite \texttt{Minted}. Cualquier comprobación fallida revierte los
	cambios realizados dentro de la misma transacción, de modo que con
	\lstinline|mintRemoteWithHeader| el anclaje no persiste si la acuñación falla. Un anclaje
	registrado en una transacción anterior, en cambio, se conserva aunque falle una solicitud de
	acuñación posterior.
\end{enumerate}

\subsection{Condiciones de aceptación y modelo de confianza}
\label{sec:puente_aceptacion}

\texttt{BeaconLightClient} solo ancla un bloque cuando el contrato verificador acepta la prueba de
consenso para la clave del programa correspondiente, la prueba declara válida la firma, condición
que el programa solo establece si la firma agregada y la rama de Merkle son correctas y
participaron en ella más de dos tercios de los miembros del comité, y el dominio, la posición del
hash y el compromiso del comité coinciden con los fijados por el propietario. \texttt{ZKReceiver}, por su parte, solo acuña cuando el verificador acepta la
prueba de inclusión para la clave de su programa, la prueba declara verificada la inclusión, la
raíz de las transacciones proviene del encabezado y la transacción pudo interpretarse, el hash de
ese encabezado está anclado, la transacción se dirigió al contrato \texttt{Sender}, los datos de
la llamada coinciden en longitud y en hash con los comprometidos por la prueba y corresponden a
una llamada a \texttt{lock} con una dirección bien formada y un monto positivo, y el par formado
por el hash del bloque y la posición de la transacción no se procesó antes.

Estas condiciones respaldan algunas de las afirmaciones que distingue la
Sección~\ref{sec:circom_limitaciones}, pero no todas. La inclusión de la transacción y la
vinculación de la raíz con un encabezado concreto quedan establecidas por la prueba de inclusión,
y el destinatario y el monto se extraen de los datos de la transacción probada, por lo que el
\textit{relayer} no puede alterarlos. La reutilización de una operación se impide mediante el
registro de los pares ya procesados. La autenticidad del encabezado, en cambio, depende de la
honestidad de la mayoría del comité de sincronización y de que el propietario de
\texttt{BeaconLightClient} mantenga actualizado el compromiso del comité, que se renueva cada 256
épocas~\cite{ethconsensusaltair}. Según la documentación de la revisión examinada, la rotación del
comité se actualiza mediante una operación del propietario y no mediante una prueba, el contrato
ancla cualquier bloque con una firma válida aunque no esté finalizado, y la prueba de inclusión no
establece que la transacción de bloqueo se haya ejecutado con éxito, para lo cual se requeriría
probar su recibo~\cite{zkguaranibridgesp1}. A ello se suman la corrección del contrato verificador
y los supuestos del sistema de prueba que lo sustenta, incluida la configuración confiable cuando
existe (Sección~\ref{sec:snarks_starks}).

El papel del \textit{relayer} admite, en consecuencia, varias distinciones. Los contratos de
destino no le otorgan un permiso exclusivo, ya que \lstinline|updateHeader|,
\lstinline|mintRemote| y \lstinline|mintRemoteWithHeader| pueden ser invocadas por cualquier
parte. Las comprobaciones de \texttt{ZKReceiver} y de los programas probados impiden, por su parte,
que quien presenta las pruebas determine el bloque, la transacción, el destinatario o el monto al
margen de lo que estas comprometen, o que reutilice una operación ya procesada. Esas comprobaciones
no bastan, sin embargo, para afirmar que toda acuñación responde a un bloqueo efectivo de tokens,
ya que el receptor no verifica el estado de ejecución de la transacción, la emisión del evento
\texttt{Locked} ni la finalidad del bloque. La detección de ese evento por el \textit{relayer}
habitual es una condición de su propio funcionamiento y no una exigencia del receptor, de modo que
las comprobaciones descritas no distinguen una llamada a \texttt{lock} incluida pero revertida, ni
un bloque firmado que no llegue a finalizarse, de una operación completada en un bloque final.
Finalmente, el \textit{relayer} sigue siendo relevante para la disponibilidad del servicio, puesto
que si deja de operar ninguna transferencia se completa hasta que otra parte genere y envíe las
pruebas. El tratamiento de las limitaciones señaladas en la versión evaluada
se discute junto con los resultados del capítulo.
